\clearpage
\section{Full translation renaming through box permutation}
\label{sec:copying-full-renaming}
\termdef{copying-full-renaming}
This milestone closes the structural bijective-renaming induction for every
constructor of the original term language, including arbitrarily nested lambda
boxes and operands at unequal allowances. The theorem is about the actual
Figure~16 translation, with its original distributors, shuffle and boxes. Its
conclusion is a finite derivation generated by the box-permutation rule and
structural congruence, followed by an ordered boundary isomorphism with a proved
resource action. It does not assert observational equivalence.

This distinction is mathematically necessary. An ordered isomorphism of two
single-box graphs fixes each outer output individually. A nontrivial permutation
of their free resources therefore requires the structural box-permutation
pre-template of Figure~21, not an enlargement of the existing isomorphism notion.
The new derivation type records this use explicitly. Its semantic interpretation
still requires the copying machine, BPerm robustness and execution congruence.
No machine or robustness hypothesis is hidden in the structural theorem.


\Needspace{12\baselineskip}
\subsection{Typed permutations and the bound-prefix offset}

\termdef{copying-box-permutation-pair}
A typed permutation of an output list is a bijection of its positions preserving
each wire label. Identity is immediate. Inversion follows by applying the forward
typing equation at the inverse position; composition follows by transitivity of
the two typing equations. The underlying finite-position bijections remain
explicit rather than being selected by choice.

\lean{Copying/BoxPermutationPair}{9}{11}{BoxPermutationPair.TypedPermutation}

\lean{Copying/BoxPermutationPair}{15}{15}{BoxPermutationPair.refl}

\lean{Copying/BoxPermutationPair}{16}{18}{BoxPermutationPair.symm}

\lean{Copying/BoxPermutationPair}{20}{22}{BoxPermutationPair.trans}

Prefix extension acts as the identity below $n$ and sends position $n+j$
to $n+p(j)$. The bound and free evaluation lemmas follow by the two branches of
this definition. For bijectivity, split a position at $n$ and use the original
inverse laws in the free branch. For typing, the prefix is homogeneous in term
wires and the suffix equation is exactly the supplied typing equation. Empty
prefixes and empty free interfaces require no special assumptions.

\lean{Copying/BoxPermutationPair}{25}{28}{BoxPermutationPair.prefixMap}

\lean{Copying/BoxPermutationPair}{30}{32}{BoxPermutationPair.prefixMap_bound}

\lean{Copying/BoxPermutationPair}{34}{36}{BoxPermutationPair.prefixMap_free}

\lean{Copying/BoxPermutationPair}{39}{62}{BoxPermutationPair.extend}

An equality cast changes only the interface index; its numerical-position
lemma follows by equality elimination. The combined term/atom permutation is
read from the already proved distributor boundary action. The two block lemmas
extract its action on term and atom indices from the corresponding port equations.

\lean{Copying/BoxPermutationPair}{65}{65}{BoxPermutationPair.cast}

\lean{Copying/BoxPermutationPair}{68}{70}{BoxPermutationPair.environment}

\lean{Copying/BoxPermutationPair}{72}{76}{BoxPermutationPair.environment_term}

\lean{Copying/BoxPermutationPair}{78}{82}{BoxPermutationPair.environment_atom}

\lean{Copying/BoxPermutationPair}{83}{85}{BoxPermutationPair.cast_val}

Splitting the body interface exposes $n$ bound variables followed by the
whole free interface. The inner permutation is the prefix extension transported
along that equality. A box-permutation pair records the outer and inner typed
bijections, pointwise fixation of bound outputs, and the equation
\[
 p_{\mathrm{body}}(n+j)=n+p_{\mathrm{free}}(j).
\]
The bound/free index injections retain the offset explicitly. Evaluation proves
that every typed outer permutation gives a pair. Conversely, splitting an
arbitrary body index at $n$ proves that a pair's inner action is uniquely
specified by its outer action. Inverse pairs follow by applying the inverse body
map to the two original equations. Composite pairs follow by substituting the
first pair's equations into the second pair's equations.

\lean{Copying/BoxPermutationPair}{90}{92}{BoxPermutationPair.interface_split}

\lean{Copying/BoxPermutationPair}{95}{97}{BoxPermutationPair.inner}

\lean{Copying/BoxPermutationPair}{99}{100}{BoxPermutationPair.boundIndex}

\lean{Copying/BoxPermutationPair}{101}{102}{BoxPermutationPair.freeIndex}

\lean{Copying/BoxPermutationPair}{105}{109}{BoxPermutationPair.Pair}

\lean{Copying/BoxPermutationPair}{111}{116}{BoxPermutationPair.inner_bound}

\lean{Copying/BoxPermutationPair}{118}{123}{BoxPermutationPair.inner_free}

\lean{Copying/BoxPermutationPair}{125}{126}{BoxPermutationPair.pair}

\lean{Copying/BoxPermutationPair}{129}{137}{BoxPermutationPair.Pair.unique}

\lean{Copying/BoxPermutationPair}{139}{148}{BoxPermutationPair.Pair.symm}

\lean{Copying/BoxPermutationPair}{150}{154}{BoxPermutationPair.Pair.trans}

\paragraph{Printed indexing discrepancy.}
\href{https://arxiv.org/pdf/1907.01257v5\#page=45}{PDF Definition~12.7(2)}
(\href{https://arxiv.org/html/1907.01257v5\#S12.SS2}{HTML Section~12.2})
prints $\rho(i)=\rho'(i-n)$ for free positions. Taken literally, its right side
omits the bound-prefix offset. With one bound and one free variable, both body
positions would map to the bound position, contradicting injectivity. The checked
regression below certifies this obstruction. We use the offset-preserving equation
above, equivalently $\rho(i)=n+\rho'(i-n)$ in the paper's one-based coordinates.
This explicit indexing correction changes neither the terms nor their translation.


\Needspace{12\baselineskip}
\subsection{Actual interface permutation and compensation}

\termdef{copying-interface-permutation}
The interface operation fixes internal ports and input ports and moves an old
output $i$ to output $p(i)$. Its inverse uses $p^{-1}$. It preserves port roles
by case analysis and wire types by the permutation's typing equation. In the
paper's convention, which reorders the boundary list, this forward-position
operation is $\Pi^{p^{-1}}_{\mathrm{id}}$. Keeping track of the inverse avoids
reversing the compensation outside a box.

\lean{Copying/InterfacePermutation}{10}{20}{InterfacePermutation.ports}

\lean{Copying/InterfacePermutation}{22}{22}{InterfacePermutation.input_iff}

\lean{Copying/InterfacePermutation}{23}{23}{InterfacePermutation.output_iff}

\lean{Copying/InterfacePermutation}{25}{29}{InterfacePermutation.label}

The permuted graph has exactly the same internal carrier, edge carrier,
finite enumerations and labels. Both incidence lists are mapped by the port
bijection in their existing order. Injectivity preserves no-duplicates and
incidence uniqueness. Pulling a port back along the inverse proves coverage;
the type equation proves the two typed-incidence laws. The resulting boundary
isomorphism is the port bijection with identity on edges, and its input and output
actions compute respectively to identity and $p$. Primitive-label restrictions
are unchanged definitionally; binding-freedom and stability follow from the
previously proved boundary-isomorphism invariance theorems.

\lean{Copying/InterfacePermutation}{33}{87}{InterfacePermutation.net}

\lean{Copying/InterfacePermutation}{90}{98}{InterfacePermutation.iso}

\lean{Copying/InterfacePermutation}{100}{100}{InterfacePermutation.iso_input}

\lean{Copying/InterfacePermutation}{101}{101}{InterfacePermutation.iso_output}

\lean{Copying/InterfacePermutation}{103}{103}{InterfacePermutation.allLabels_iff}

\lean{Copying/InterfacePermutation}{105}{105}{InterfacePermutation.bindingFree_iff}

\lean{Copying/InterfacePermutation}{107}{108}{InterfacePermutation.stable_iff}

\termdef{copying-boundary-compensation}
Suppose $f:G\to H$ fixes inputs and has output action $p$. Compose it with
interface permutation by $p^{-1}$ on $H$. The composite fixes both boundaries,
so the existing conversion theorem yields a genuine boundary-fixing recursive
isomorphism $G\cong \Pi^p_{\mathrm{id}}H$. Applying this to the interface
operation itself proves the round-trip law. Applying the same permutation to
both sides of a fixed-boundary isomorphism follows by conjugating its port map.
Finally, permuting the compensated isomorphism and cancelling the round trip
factors $f$ through actual interface permutation of $G$. All these results preserve
the original ordered notion of recursive isomorphism.

\lean{Copying/InterfacePermutation}{111}{132}{InterfacePermutation.compensate}

\lean{Copying/InterfacePermutation}{135}{136}{InterfacePermutation.roundtrip}

\lean{Copying/InterfacePermutation}{139}{153}{InterfacePermutation.congr}

\lean{Copying/BoxPermutationRule}{10}{13}{BoxPermutationRule.realize}

\Needspace{12\baselineskip}
\subsection{The concrete BPerm graph pair and its structural safety}

For a body $G$ and a certified pair $(p_{\mathrm{body}},p_{\mathrm{free}})$,
the literal Figure~21 rule relates a box containing
$\Pi^{p_{\mathrm{body}}}_{\mathrm{id}}G$ to
$\Pi^{p_{\mathrm{free}}}_{\mathrm{id}}(\operatorname{box}G)$.
In our forward-position convention both displayed permutations therefore use
inverse bijections. The raw relation has exactly this constructor, with the
same outer interfaces. Its isomorphism closure allows a boundary-fixing recursive
isomorphism before and after that one raw rule.

A useful compensated normal form starts from the original box and ends at the
box of the forward-permuted body, with inverse free permutation outside. The
function \texttt{right} constructs this graph; the normal-form theorem below
\emph{derives} the pair from the literal rule. Neither relation asserts
observational equivalence or execution.

\lean{Copying/BoxPermutationRule}{22}{24}{BoxPermutationRule.right}

\lean{Copying/BoxPermutationRule}{28}{31}{BoxPermutationRule.Step}

\lean{Copying/BoxPermutationRule}{34}{36}{BoxPermutationRule.UpToIso}

Every recursively checked primitive-label restriction is preserved because
the only new outer label is the same box and the body's labels have not changed.
If a body isomorphism has the pair's prescribed output action, compensation
identifies its source with the inverse-permuted target. Apply box congruence to
that fixed-boundary isomorphism, then apply the literal paper rule to the target
body. This places the required pair in the rule's isomorphism closure without
assuming the desired rule instance. Specialising to the actual interface-operation
isomorphism derives the compensated normal form.

\lean{Copying/BoxPermutationRule}{40}{45}{BoxPermutationRule.allLabels_right}

\lean{Copying/BoxPermutationRule}{48}{54}{BoxPermutationRule.of_bodyIso}

\lean{Copying/BoxPermutationRule}{57}{58}{BoxPermutationRule.normalized}

For comparison, take an ordered isomorphism between two original box
corollas. Apply its target-incidence equation to position $i$ in the ordered
output list. The result says that output $i$ maps to output $i$. Hence a
nontrivial free permutation cannot be implemented by such an isomorphism,
regardless of the two body graphs. This rules out the tempting but incorrect
shortcut of treating BPerm as ordinary graph renaming.

\lean{Copying/BoxPermutationRule}{61}{67}{BoxPermutationRule.box_output_fixed}

\lean{Copying/BoxPermutationRule}{70}{78}{BoxPermutationRule.nontrivial_not_ordered}

Binding-freedom of a single box is equivalent to binding-freedom of its
body: the outer level contains no hole, while focus exclusion and inner binding
paths are exactly those of the body. Combine this with interface-permutation
invariance to obtain the right graph's safety equivalence. Case analysis on a
raw rule proves preservation of all primitive restrictions and binding-freedom.
The same properties extend across the two boundary-fixing isomorphisms.
These are structural safety theorems, not machine input-safety or robustness.

\lean{Copying/BoxPermutationSafety}{10}{18}{BoxPermutationSafety.bindingFree_box_iff}

\lean{Copying/BoxPermutationSafety}{21}{23}{BoxPermutationSafety.bindingFree_right}

\lean{Copying/BoxPermutationSafety}{25}{31}{BoxPermutationSafety.Step.allLabels}

\lean{Copying/BoxPermutationSafety}{33}{37}{BoxPermutationSafety.Step.bindingFree}

\lean{Copying/BoxPermutationSafety}{40}{44}{BoxPermutationSafety.UpToIso.bindingFree}

\lean{Copying/BoxPermutationSafety}{46}{51}{BoxPermutationSafety.UpToIso.allLabels}

\Needspace{12\baselineskip}
\subsection{The original lambda interface and abstraction shell}

Lift a variable bijection under a binder using the existing capture-free
\texttt{Term.lift}: zero is fixed and successor positions follow the free
bijection. The inverse equations follow by zero/successor case analysis. Combine
this lifted bijection with the atom bijection and apply precisely the interface
cast already used by \texttt{lambdaBox}. The resulting lambda pair satisfies the
bound/free conditions: split free positions into term and atom blocks and use
their exact distributor actions. This proves the agreement between the scoped
binder convention and the box's ordered body boundary.

\lean{Copying/LambdaPermutationPair}{10}{14}{LambdaPermutationPair.lift}

\lean{Copying/LambdaPermutationPair}{17}{20}{LambdaPermutationPair.lambdaBodyPermutation}

\lean{Copying/LambdaPermutationPair}{22}{45}{LambdaPermutationPair.lambdaPair}

A resource-action certificate determines the whole output permutation:
split an arbitrary position into its term or atom block and apply the respective
certificate field. Equality transport carries this equation through the lambda
body cast. The general body-isomorphism rule then produces a BPerm step for the
original lambda box, with the inverse outer permutation explicitly present.

\lean{Copying/LambdaPermutationPair}{50}{66}{LambdaPermutationPair.outputs}

\lean{Copying/LambdaPermutationPair}{71}{75}{LambdaPermutationPair.cast_output_action}

\lean{Copying/LambdaPermutationPair}{81}{87}{LambdaPermutationPair.lambda_of_action}

The abstraction shell is the literal lambda operation, followed by a
supplied box graph, followed by the mandatory single-use distributors. Its
singleton thunk input fixes compatibility with the lambda edge. Build the prefix
isomorphism; use its typed free-output action to choose the existing single-use
forest automorphism. The matching equation proves serial compatibility. Composing
the two isomorphisms gives shell congruence, and the final port equation proves
that the shell retains exactly the supplied box's resource action.

\lean{Copying/AbstractionShell}{11}{13}{AbstractionShell.graph}

\lean{Copying/AbstractionShell}{16}{29}{AbstractionShell.prefix_compatible}

\lean{Copying/AbstractionShell}{31}{34}{AbstractionShell.prefixIso}

\lean{Copying/AbstractionShell}{37}{39}{AbstractionShell.forestIso}

\lean{Copying/AbstractionShell}{41}{48}{AbstractionShell.forest_compatible}

\lean{Copying/AbstractionShell}{51}{52}{AbstractionShell.iso}

\lean{Copying/AbstractionShell}{55}{60}{AbstractionShell.iso_output}

\Needspace{12\baselineskip}
\subsection{Finite structural derivations and their precise scope}

\termdef{copying-bperm-derivation}
A derivation is an inductively generated finite proof tree. Its generators are
boundary-fixing recursive isomorphism and the concrete BPerm rule. It is closed
under symmetry, transitivity, serial and parallel composition, tensor, enclosing
a box, equality transport of interfaces, and structural depth raising. These are
explicit syntactic congruence constructors. In particular, the raising constructor
is not a claim of execution-padding invariance, and the rule constructor is not
a proof of BPerm robustness. There is no constructor for arbitrary pairs of terms,
no constructor assuming renaming, and no Proposition~12.2 constructor.

\lean{Copying/BoxPermutationDerivation}{12}{28}{BoxPermutationDerivation.Derivation}

Reflexivity uses identity isomorphism. A rule up to isomorphism expands
into two isomorphism generators and one raw rule. Depth equality is eliminated;
finite padding follows by induction on its number of raises, and common-allowance
padding adds the existing equality cast. Shell, abstraction and unary congruence
are derived by serial composition, retaining each fixed component. Pair congruence
uses the original parallel boundary and its original output-interface cast;
application and binary congruence prefix the original operation and append the
original two-use distributors. Thus the derived clauses preserve Figure~16
literally.

\lean{Copying/BoxPermutationDerivation}{32}{32}{BoxPermutationDerivation.refl}

\lean{Copying/BoxPermutationDerivation}{34}{36}{BoxPermutationDerivation.of_upToIso}

\lean{Copying/BoxPermutationDerivation}{38}{39}{BoxPermutationDerivation.depthCast}

\lean{Copying/BoxPermutationDerivation}{41}{45}{BoxPermutationDerivation.pad}

\lean{Copying/BoxPermutationDerivation}{47}{49}{BoxPermutationDerivation.padTo}

\lean{Copying/BoxPermutationDerivation}{51}{53}{BoxPermutationDerivation.shell}

\lean{Copying/BoxPermutationDerivation}{55}{57}{BoxPermutationDerivation.abstraction}

\lean{Copying/BoxPermutationDerivation}{59}{61}{BoxPermutationDerivation.unary}

\lean{Copying/BoxPermutationDerivation}{62}{70}{BoxPermutationDerivation.pair}

\lean{Copying/BoxPermutationDerivation}{72}{75}{BoxPermutationDerivation.application}

\lean{Copying/BoxPermutationDerivation}{77}{80}{BoxPermutationDerivation.binary}

\Needspace{12\baselineskip}
\subsection{Full syntax induction with a resource-action invariant}

\termdef{copying-renaming-witness}
A renaming witness consists of an intermediate graph $M$, a finite derivation
$G\Rightarrow M$, an ordered boundary isomorphism $M\to H$, and its exact
variable/atom resource-action certificate. The datatype permits heterogeneous
allowance indices at its boundary, but its constructor requires those indices to
coincide; consequently every witness proves equality of the allowances. The
intermediate graph is constructed in the proof, not supplied by the caller.
An ordinary certified boundary isomorphism gives a witness with a reflexive
derivation.

\lean{Copying/FullTranslationRenaming}{13}{17}{FullTranslationRenaming.Renaming}

\lean{Copying/FullTranslationRenaming}{22}{24}{FullTranslationRenaming.of_iso}

\lean{Copying/FullTranslationRenaming}{26}{27}{FullTranslationRenaming.depth_eq}

Unary congruence combines the derived syntactic clause with the existing
unary isomorphism and its action law. Padding combines derivation padding with
proved isomorphism padding and certificate preservation. For application and
binary operations, first match the two resource actions using their certificates,
then use the existing constructor isomorphism. The padded variants eliminate the
allowance equalities supplied by the two witnesses and pad both operands to
exactly the original maximum. Thus neither equal original operand depths nor a
separate matching premise is needed.

\lean{Copying/FullTranslationRenaming}{29}{33}{FullTranslationRenaming.unary}

\lean{Copying/FullTranslationRenaming}{35}{40}{FullTranslationRenaming.padTo}

\lean{Copying/FullTranslationRenaming}{42}{52}{FullTranslationRenaming.application}

\lean{Copying/FullTranslationRenaming}{54}{64}{FullTranslationRenaming.binary}

\lean{Copying/FullTranslationRenaming}{66}{77}{FullTranslationRenaming.padded_application}

\lean{Copying/FullTranslationRenaming}{79}{90}{FullTranslationRenaming.padded_binary}

The abstraction case is the essential new step. Lift the body's existing
derivation through the original abstraction constructor. Its remaining body
isomorphism supplies the lambda BPerm instance proved above. Lift that single
rule, with its two isomorphisms, through the abstraction shell. The resulting
box is the target box with inverse free-interface permutation. Invert that
interface operation and use shell congruence, which also moves the original
distributors. Its output equation gives the required free variable/atom action.
The new witness therefore has exactly the required shape for another enclosing
lambda; arbitrary nesting follows without a new assumption.

\lean{Copying/FullTranslationRenaming}{93}{111}{FullTranslationRenaming.abstraction}

The final isomorphism of any witness can be compensated, yielding a finite
derivation from $G$ to the explicitly boundary-compensated target. The public
translation theorem is proved by induction on the original term: leaves use the
certified atomic isomorphisms; unary, application and binary nodes use the
constructor results; lambda uses the capture-free lifted variable bijection and
the new abstraction result. All arities and all finite variable/atom environments
are quantified. The conclusion names the original translated term and the
translation of its original syntactic renaming, with no replacement translation.

\lean{Copying/FullTranslationRenaming}{115}{121}{FullTranslationRenaming.derivation}

\lean{Copying/FullTranslationRenaming}{127}{137}{FullTranslationRenaming.term_rename_via_boxPermutation}

\Needspace{12\baselineskip}
\subsection{Invariants, regressions and remaining semantic obligations}

Serial expansion contains exactly its two components' labels; inspect the
Boolean component tag in each direction. Induction on structural derivations
then proves preservation of every recursively checked primitive-label predicate.
Isomorphisms and raw rules use their existing invariance results; composition
uses component conjunctions; parallel assembly uses its component tag; boxing
uses the body predicate; casts eliminate equalities; raising uses structural
padding invariance. Focus exclusion, hole exclusion and membership in the full
example language are immediate specialisations.

\lean{Copying/BoxPermutationInvariants}{9}{14}{BoxPermutationInvariants.allLabels_iff}

\lean{Copying/BoxPermutationInvariants}{21}{42}{BoxPermutationInvariants.allLabels}

\lean{Copying/BoxPermutationInvariants}{44}{44}{BoxPermutationInvariants.focusFree}

\lean{Copying/BoxPermutationInvariants}{45}{45}{BoxPermutationInvariants.holeFree}

\lean{Copying/BoxPermutationInvariants}{46}{46}{BoxPermutationInvariants.language}

The public derivation corollary exposes the final allowance equality and
compensated target. Two further corollaries apply the full theorem to the exact
left and right programs of Proposition~12.2; they prove renaming of each program,
not equivalence between those two programs. The nested regression uses two binders
above an application and nontrivially swaps both free resource sorts. A constant
beside that term tests unequal allowances. The named-environment corollary simply
uses the original distinct-list environment wrapper.

\lean{Copying/FullRenamingExamples}{9}{15}{FullRenamingExamples.termRenamingDerivation}

\lean{Copying/FullRenamingExamples}{18}{21}{FullRenamingExamples.privateCell_fullRenaming}

\lean{Copying/FullRenamingExamples}{23}{26}{FullRenamingExamples.constantResult_fullRenaming}

\lean{Copying/FullRenamingExamples}{29}{29}{FullRenamingExamples.nestedRenamingTerm}

\lean{Copying/FullRenamingExamples}{31}{34}{FullRenamingExamples.nestedRenaming}

\lean{Copying/FullRenamingExamples}{37}{41}{FullRenamingExamples.unequalDepthRenaming}

\lean{Copying/FullRenamingExamples}{44}{50}{FullRenamingExamples.full_namedEnvironment}

The negative regressions check three separate errors. Swapping a bound
position with a free position contradicts the bound-fixation condition. Swapping
two free outputs of an ordered box contradicts the corolla output theorem.
Finally, the printed unshifted pair equation at one bound and one free variable
forces a collision in a two-element bijection, exposing the indexing discrepancy
rather than silently repairing it.

\lean{Copying/FullRenamingExamples}{53}{63}{FullRenamingExamples.captured_swap_rejected}

\lean{Copying/FullRenamingExamples}{66}{76}{FullRenamingExamples.box_swap_not_iso}

\lean{Copying/FullRenamingExamples}{80}{87}{FullRenamingExamples.unshifted_box_pair_impossible}

\paragraph{Completion boundary and estimate.}
As of 26 October 2026, the structural renaming induction covers all syntax,
including lambda boxes. The raw BPerm pair and its structural safety laws are
also complete. Interpreting these derivations as contextual equivalences remains
open: it requires BPerm robustness, the concrete copying UAM and the corresponding
execution and congruence laws. Non-bijective graph renaming, graph substitution,
the other template families, semantic weakening and Figure~22's final chain also
remain. Appendix~\ref{app:remaining-work} lists only unfinished obligations.
The estimated mathematical completion is now roughly \textbf{50--55\%}, with a
broad plausible range of \textbf{35--65\%}. This is an effort estimate, not a count
of declarations or a claim that half of the machine rules have been implemented.
Proposition~12.2 is still unproved.
